% !TeX TS-program = xelatex

\documentclass[dvipsnames]{beamer}
\usetheme{metropolis}

\usepackage{multicol}

\usepackage{ragged2e} % who justifies the text
\usepackage{xecolor}
\usepackage{amsmath}
%\usefonttheme[onlymath]{serif} % change the math font

\usepackage{tabularx}
\usepackage{booktabs}
\usepackage[style=numeric,sorting=ynt]{biblatex}
\addbibresource{references.bib}
\usepackage[localise]{xepersian}

\settextfont{Vazir}
\setlatintextfont{Roboto}

%---------------------------------------------------------------------------------
% Settings to force Beamer works with Xepersian and RTL typesetting
%---------------------------------------------------------------------------------
%\raggedleft

% For right to left lists (itemize and enumerate)
\makeatletter
\newcommand{\لیست‌فارسی}{\raggedleft\rightskip\@totalleftmargin}
\makeatother
% Correct the bullet for RTL texts
\setbeamertemplate{itemize item}{\scriptsize\raise1.25pt%
 \hbox{\donotcoloroutermaths$\blacktriangleleft$}} 

% To force beamer use numbering in captions
\setbeamertemplate{caption}[numbered]{}% Number float-like environments

\setbeamertemplate{footline}[frame number]
\setbeamertemplate{section in toc}[circle]
\setbeamertemplate{blocks}[rounded][shadow=true]
\setbeamercolor{block body}{bg=lightgray}
\setbeamercolor{headline}{bg=orange}
\setbeamersize{text margin left=1cm,text margin right=1cm}

\setbeamertemplate{headline}
{
		\begin{beamercolorbox}{section in head/foot}
				\vspace{2pt}\insertnavigation{\paperwidth}\vspace{2pt}
		\end{beamercolorbox}
}

%---------------------------------------------------------------------------------
% To force beamer use numbering in captions
\setbeamertemplate{caption}[numbered]{}% Number float-like environments

\setbeamertemplate{footline}
{%
	\leavevmode%
	\hbox{%
		\begin{beamercolorbox}[wd=.333333\paperwidth,ht=2.25ex,dp=1ex,center]{author in head/foot}%
			\usebeamerfont{author in head/foot}\insertshortauthor%
		\end{beamercolorbox}%
		\begin{beamercolorbox}[wd=.333333\paperwidth,ht=2.25ex,dp=1ex,center]{title in head/foot}%
			\usebeamerfont{title in head/foot}\insertshorttitle%
		\end{beamercolorbox}%
		\begin{beamercolorbox}[wd=.333333\paperwidth,ht=2.25ex,dp=1ex,right]{date in head/foot}%
			\usebeamerfont{date in head/foot}\insertsection\hspace*{2em}
			\insertframenumber~ از \inserttotalframenumber{} \hspace*{2ex}%
		\end{beamercolorbox}
	}%
}
\setbeamertemplate{section in toc}[circle]
\setbeamertemplate{blocks}[rounded][shadow=true]
\setbeamercolor{block title}{bg=orange}
\setbeamercolor{block body}{bg=lightgray}
\setbeamersize{text margin left=1cm,text margin right=1cm}
\setbeamertemplate{frametitle continuation}{\insertcontinuationcount}

%---------------------------------------------------------------------------------
\title{ارزیابی کارآیی شبکه‌های توان پایین}
\subtitle{}
\author{پرهام الوانی}
\institute{%
		دانشکده مهندسی کامپیوتر\\
		دکتر مهدی راستی
}
\date{\today}
\titlegraphic{\vspace{4.5cm}\flushleft\includegraphics[height=50pt]{images/logo}}

\begin{document}

\makeatletter

\setbeamertemplate{title}{%
	\linespread{1.0}%
	\inserttitle%
	\par%
	\vspace*{0.5em}
}
\setbeamertemplate{subtitle}{%
	\insertsubtitle%
	\par%
	\vspace*{0.5em}
}

\AtBeginSection[]
{%
	\begin{frame}{فهرست}
		\tableofcontents[currentsection]
	\end{frame}
	\begin{frame}
		\begin{center}
			\insertsectionnumber. \insertsection%
		\end{center}
		\usebeamertemplate*{title separator}
	\end{frame}
}

\makeatother


\begin{persian}
	%------------------------------------------
	% Title frame (0)
	%------------------------------------------
	{%
		\setbeamertemplate{footline}{}
		\begin{frame}
			\titlepage%
		\end{frame}
	}

	%-------------------------------------------------------------------------------
	\begin{frame}{فهرست}
		\tableofcontents[pausesections]
	\end{frame}

	%-------------------------------------------------------------------------------
	\section{شبکه‌های \متن‌لاتین{LoRa}}

	\section{کارهای مرتبط}
	\begin{frame}{}
		\begin{table}[h]
			\caption{مرور پژوهش‌های حوزه ارزیابی کارایی شبکه‌های توان پایین در لایه دسترسی}
			\vspace{0.3cm}
			\tiny{}
			\begin{tabularx}{\textwidth}{XXXXXXXX}
				\toprule
				مرجع &
				تحرک &
				\متن‌لاتین{coverage} &
				شبکه &
				لایه انتقال &
				لایه اپلیکشن &
				اندازه بسته &
				\متن‌لاتین{Coding Rate} \\
				\midrule
				\midrule
				\cite{sensors-18-00772-v3} &
				سیار &
				گزارش شده &
				\متن‌لاتین{LoRaWAN} &
				ندارد &
				ندارد &
				گزارش شده &
			  	گزارش شده \\
				\midrule
				\bottomrule
			\end{tabularx}
		\end{table}
	\end{frame}


	\begin{frame}{\cite{sensors-18-00772-v3}}
	  \شروع{فقرات}
	  \لیست‌فارسی
	  \فقره یک نود متحرک که پارامترهای ارسالش با زمان تغییر می‌کند.
	  \فقره ناحیه‌های شهری، نیمه‌شهری و غیرشهری
	  \فقره ارزیابی رادیویی بر پایه مدل انتشار \متن‌لاتین{Okumura-Hata} پیش از انجام ارزیابی عملیاتی
	  \پایان{فقرات}
	\end{frame}

	%-------------------------------------------------------------------------------
	\begin{frame}[allowframebreaks]{مراجع}
		\begin{latin}
		\printbibliography[title=مراجع]
		\end{latin}
	\end{frame}

\end{persian}
\end{document}
